% =====================================================================
%  BUỔI 19 — ĐÁP ÁN ĐỀ CHỮA (từ hoc-sinh/02-dap-an.md)
%  Sinh tự động bằng scripts/md2tex-chuade.py — không sửa tay.
% =====================================================================
\documentclass[11pt, a4paper]{article}
% =====================================================================
%  CAE SHANGHAI — GÓI LỆNH DÙNG CHUNG CHO TÀI LIỆU LUYỆN THI CSCA TOÁN
%  Biên dịch bằng XeLaTeX. Mọi file .tex trong buổi học \input file này.
% =====================================================================

% ---------- NGÔN NGỮ & PHÔNG CHỮ ----------
\usepackage{fontspec}
\usepackage{polyglossia}
\setmainlanguage{vietnamese}
\setotherlanguage{english}

% Chữ Latin (Anh + Việt): TeX Gyre Termes — bản clone Times New Roman,
% đi kèm TeX Live, hỗ trợ đầy đủ dấu tiếng Việt.
\setmainfont{TeX Gyre Termes}
\setsansfont{TeX Gyre Heros}

% Chữ Hán: ưu tiên Source Han Serif SC (思源宋体) — chuẩn học thuật Trung Quốc.
% Nếu máy chưa cài, tự động dùng Songti SC (宋体, có sẵn trên macOS).
% Cả hai đều thuộc họ 宋体 nên diện mạo tài liệu không đổi.
\IfFontExistsTF{Source Han Serif SC}{
  \newcommand{\cjkfamilyname}{Source Han Serif SC}
}{
  \IfFontExistsTF{Songti SC}{
    \newcommand{\cjkfamilyname}{Songti SC}
  }{
    \newcommand{\cjkfamilyname}{Noto Serif CJK SC}
  }
}

% xeCJK tự chuyển font khi gặp ký tự Hán, không cần bọc \cjk{} thủ công.
% Đây là điều kiện bắt buộc: font Latin (TeX Gyre Termes) không có glyph chữ Hán,
% nếu không khai báo xeCJK thì mọi chữ Hán viết trần sẽ bị nuốt khỏi bản in.
\usepackage{xeCJK}
\setCJKfamilyfont{caehan}{\cjkfamilyname}[Scale=0.95]
\setCJKmainfont{\cjkfamilyname}[Scale=0.95]
\setCJKsansfont{\cjkfamilyname}[Scale=0.95]
\setCJKmonofont{\cjkfamilyname}[Scale=0.95]

% \cjk{} ép dùng font Hán trong ngữ cảnh cần thiết (ví dụ trong tiêu đề in đậm
% hoặc trong bảng có định dạng riêng). Dùng \CJKfamily chứ không phải tên font,
% vì tên font chỉ là văn bản — sẽ in ra chữ thay vì đổi font.
\newcommand{\cjk}[1]{{\CJKfamily{caehan}#1}}

% ---------- TOÁN HỌC ----------
\usepackage{amsmath, amssymb, amsthm, mathtools}
\usepackage{unicode-math}
\setmathfont{TeX Gyre Termes Math}

% ---------- BỐ CỤC TRANG ----------
% Lề phải chừa vùng an toàn của letterhead:
%   banner trên chiếm 0–19mm  → top tối thiểu 2,2cm
%   banner dưới chiếm 281–297mm → bottom tối thiểu 2,2cm
% Lề trên 2,6cm (26mm) để tiêu đề đầu trang không đè banner (banner hết ở 19mm).
\usepackage[a4paper, top=2.6cm, bottom=2.4cm, left=2.5cm, right=2.0cm]{geometry}
\usepackage{fancyhdr}
\usepackage{titlesec}
\usepackage{enumitem}
\usepackage{graphicx}
\usepackage{booktabs, array, longtable, tabularx}
\usepackage{xcolor}
\usepackage[most]{tcolorbox}
\usepackage[colorlinks=true, linkcolor=cae_blue, citecolor=cae_blue, urlcolor=cae_blue]{hyperref}

% ---------- BẢNG MÀU THƯƠNG HIỆU ----------
\definecolor{cae_blue}{RGB}{0, 51, 102}
\definecolor{cae_gold}{RGB}{196, 149, 60}
\definecolor{box_bg}{RGB}{245, 247, 250}
\definecolor{box_line}{RGB}{0, 51, 102}

% ---------- TIÊU ĐỀ MỤC ----------
% Đệm đầu trang: banner letterhead chiếm 0–19mm, mà lề trên là 22mm.
% \section* đầu trang (sau \clearpage) sát lề nên phải chừa thêm, nếu không
% sẽ đè lên banner.
\titleformat{\section}
  {\normalfont\large\bfseries\sffamily\color{cae_blue}}{\thesection.}{0.6em}{}
\titleformat{\subsection}
  {\normalfont\normalsize\bfseries\sffamily\color{cae_blue!85!black}}{\thesubsection}{0.6em}{}
\titlespacing*{\section}{0pt}{14pt}{6pt}
\titlespacing*{\subsection}{0pt}{10pt}{4pt}

% Tiêu đề không đánh số ở đầu trang: \section* sau \clearpage nằm sát lề trên
% nên đè lên banner letterhead. Dùng \vspace* ngay TRƯỚC \section* không đủ,
% vì titlesec chèn khoảng cách riêng. Cách chắc chắn là đệm bằng \null.
\newcommand{\sectiondautrang}[1]{\section*{#1}}

% Dạng bài dùng \subsection*{} nên cần định dạng riêng, đậm và có gạch chân màu
\titleformat{\subsubsection}
  {\normalfont\normalsize\bfseries\sffamily\color{cae_blue!85!black}}{\thesubsubsection}{0.6em}{}

% ---------- LETTERHEAD (ảnh nền toàn trang) ----------
% Ảnh letterhead của CAE Shanghai: banner xanh trên/dưới kèm logo chìm giữa.
% Tỉ lệ ảnh 1414x2000 = 0.7070, khớp A4 (0.7071) nên chèn thẳng không méo.
% Đặt ở lớp nền dưới cùng để chữ không bị ảnh che.
% Dùng \AddToShipoutPictureBG (không có dấu *) để áp dụng cho MỌI trang.
\usepackage{eso-pic}
\AddToShipoutPictureBG{%
  \AtPageLowerLeft{%
    \includegraphics[width=\paperwidth,height=\paperheight]{images/letterhead.png}%
  }%
}

% ---------- ĐẦU / CHÂN TRANG ----------
% Letterhead đã có banner trên và dưới, nên bỏ dòng chạy để tránh chồng chữ.
% Số trang đặt trong vùng nội dung, sát trên banner dưới, lề phải.
\pagestyle{fancy}
\setlength{\headheight}{0pt}
\fancyhf{}
\renewcommand{\headrulewidth}{0pt}
\renewcommand{\footrulewidth}{0pt}
\fancyfoot[R]{\raisebox{0.9cm}{\small\color{cae_blue}\thepage}}

% ---------- MÔI TRƯỜNG NỘI DUNG ----------
% Trong tcolorbox, LaTeX đặt parindent = parskip = 0, nên các đoạn văn trong
% khung dính liền nhau. Đặt lại khoảng cách đoạn cho mọi khung.
\newcommand{\khoangcachdoan}{\setlength{\parskip}{0.5em}\setlength{\parindent}{0pt}}
\tcbset{before upper={\khoangcachdoan}}

\newtcolorbox{dinhnghia}[1][]{
  colback=box_bg, colframe=box_line, boxrule=0.6pt, arc=1pt,
  fonttitle=\bfseries, title={Định nghĩa #1}, breakable}
\newtcolorbox{dinhly}[1][]{
  colback=box_bg, colframe=cae_gold, boxrule=0.6pt, arc=1pt,
  fonttitle=\bfseries, title={Tính chất #1}, breakable}
\newtcolorbox{vidu}[1][]{
  colback=white, colframe=black!40, boxrule=0.4pt, arc=1pt,
  fonttitle=\bfseries, title={Ví dụ #1}, breakable}
\newtcolorbox{baitap}[1][]{
  colback=white, colframe=cae_blue!50, boxrule=0.4pt, arc=1pt,
  fonttitle=\bfseries, title={Câu #1}, breakable}

% Ô đáp án: viền xanh đậm hơn, dùng trong file đáp án
\newtcolorbox{dapan}[1][]{
  colback=box_bg, colframe=cae_blue, boxrule=0.6pt, arc=1pt,
  fonttitle=\bfseries, title={Câu #1}, breakable}

% ---------- TÊN CHỨNG MINH ----------
\renewcommand{\proofname}{\normalfont\bfseries Chứng minh}

% ---------- MÔN HỌC ----------
% Đổi được ở file .tex trước khi % =====================================================================
%  CAE SHANGHAI — GÓI LỆNH DÙNG CHUNG CHO TÀI LIỆU LUYỆN THI CSCA TOÁN
%  Biên dịch bằng XeLaTeX. Mọi file .tex trong buổi học \input file này.
% =====================================================================

% ---------- NGÔN NGỮ & PHÔNG CHỮ ----------
\usepackage{fontspec}
\usepackage{polyglossia}
\setmainlanguage{vietnamese}
\setotherlanguage{english}

% Chữ Latin (Anh + Việt): TeX Gyre Termes — bản clone Times New Roman,
% đi kèm TeX Live, hỗ trợ đầy đủ dấu tiếng Việt.
\setmainfont{TeX Gyre Termes}
\setsansfont{TeX Gyre Heros}

% Chữ Hán: ưu tiên Source Han Serif SC (思源宋体) — chuẩn học thuật Trung Quốc.
% Nếu máy chưa cài, tự động dùng Songti SC (宋体, có sẵn trên macOS).
% Cả hai đều thuộc họ 宋体 nên diện mạo tài liệu không đổi.
\IfFontExistsTF{Source Han Serif SC}{
  \newcommand{\cjkfamilyname}{Source Han Serif SC}
}{
  \IfFontExistsTF{Songti SC}{
    \newcommand{\cjkfamilyname}{Songti SC}
  }{
    \newcommand{\cjkfamilyname}{Noto Serif CJK SC}
  }
}

% xeCJK tự chuyển font khi gặp ký tự Hán, không cần bọc \cjk{} thủ công.
% Đây là điều kiện bắt buộc: font Latin (TeX Gyre Termes) không có glyph chữ Hán,
% nếu không khai báo xeCJK thì mọi chữ Hán viết trần sẽ bị nuốt khỏi bản in.
\usepackage{xeCJK}
\setCJKfamilyfont{caehan}{\cjkfamilyname}[Scale=0.95]
\setCJKmainfont{\cjkfamilyname}[Scale=0.95]
\setCJKsansfont{\cjkfamilyname}[Scale=0.95]
\setCJKmonofont{\cjkfamilyname}[Scale=0.95]

% \cjk{} ép dùng font Hán trong ngữ cảnh cần thiết (ví dụ trong tiêu đề in đậm
% hoặc trong bảng có định dạng riêng). Dùng \CJKfamily chứ không phải tên font,
% vì tên font chỉ là văn bản — sẽ in ra chữ thay vì đổi font.
\newcommand{\cjk}[1]{{\CJKfamily{caehan}#1}}

% ---------- TOÁN HỌC ----------
\usepackage{amsmath, amssymb, amsthm, mathtools}
\usepackage{unicode-math}
\setmathfont{TeX Gyre Termes Math}

% ---------- BỐ CỤC TRANG ----------
% Lề phải chừa vùng an toàn của letterhead:
%   banner trên chiếm 0–19mm  → top tối thiểu 2,2cm
%   banner dưới chiếm 281–297mm → bottom tối thiểu 2,2cm
% Lề trên 2,6cm (26mm) để tiêu đề đầu trang không đè banner (banner hết ở 19mm).
\usepackage[a4paper, top=2.6cm, bottom=2.4cm, left=2.5cm, right=2.0cm]{geometry}
\usepackage{fancyhdr}
\usepackage{titlesec}
\usepackage{enumitem}
\usepackage{graphicx}
\usepackage{booktabs, array, longtable, tabularx}
\usepackage{xcolor}
\usepackage[most]{tcolorbox}
\usepackage[colorlinks=true, linkcolor=cae_blue, citecolor=cae_blue, urlcolor=cae_blue]{hyperref}

% ---------- BẢNG MÀU THƯƠNG HIỆU ----------
\definecolor{cae_blue}{RGB}{0, 51, 102}
\definecolor{cae_gold}{RGB}{196, 149, 60}
\definecolor{box_bg}{RGB}{245, 247, 250}
\definecolor{box_line}{RGB}{0, 51, 102}

% ---------- TIÊU ĐỀ MỤC ----------
% Đệm đầu trang: banner letterhead chiếm 0–19mm, mà lề trên là 22mm.
% \section* đầu trang (sau \clearpage) sát lề nên phải chừa thêm, nếu không
% sẽ đè lên banner.
\titleformat{\section}
  {\normalfont\large\bfseries\sffamily\color{cae_blue}}{\thesection.}{0.6em}{}
\titleformat{\subsection}
  {\normalfont\normalsize\bfseries\sffamily\color{cae_blue!85!black}}{\thesubsection}{0.6em}{}
\titlespacing*{\section}{0pt}{14pt}{6pt}
\titlespacing*{\subsection}{0pt}{10pt}{4pt}

% Tiêu đề không đánh số ở đầu trang: \section* sau \clearpage nằm sát lề trên
% nên đè lên banner letterhead. Dùng \vspace* ngay TRƯỚC \section* không đủ,
% vì titlesec chèn khoảng cách riêng. Cách chắc chắn là đệm bằng \null.
\newcommand{\sectiondautrang}[1]{\section*{#1}}

% Dạng bài dùng \subsection*{} nên cần định dạng riêng, đậm và có gạch chân màu
\titleformat{\subsubsection}
  {\normalfont\normalsize\bfseries\sffamily\color{cae_blue!85!black}}{\thesubsubsection}{0.6em}{}

% ---------- LETTERHEAD (ảnh nền toàn trang) ----------
% Ảnh letterhead của CAE Shanghai: banner xanh trên/dưới kèm logo chìm giữa.
% Tỉ lệ ảnh 1414x2000 = 0.7070, khớp A4 (0.7071) nên chèn thẳng không méo.
% Đặt ở lớp nền dưới cùng để chữ không bị ảnh che.
% Dùng \AddToShipoutPictureBG (không có dấu *) để áp dụng cho MỌI trang.
\usepackage{eso-pic}
\AddToShipoutPictureBG{%
  \AtPageLowerLeft{%
    \includegraphics[width=\paperwidth,height=\paperheight]{images/letterhead.png}%
  }%
}

% ---------- ĐẦU / CHÂN TRANG ----------
% Letterhead đã có banner trên và dưới, nên bỏ dòng chạy để tránh chồng chữ.
% Số trang đặt trong vùng nội dung, sát trên banner dưới, lề phải.
\pagestyle{fancy}
\setlength{\headheight}{0pt}
\fancyhf{}
\renewcommand{\headrulewidth}{0pt}
\renewcommand{\footrulewidth}{0pt}
\fancyfoot[R]{\raisebox{0.9cm}{\small\color{cae_blue}\thepage}}

% ---------- MÔI TRƯỜNG NỘI DUNG ----------
% Trong tcolorbox, LaTeX đặt parindent = parskip = 0, nên các đoạn văn trong
% khung dính liền nhau. Đặt lại khoảng cách đoạn cho mọi khung.
\newcommand{\khoangcachdoan}{\setlength{\parskip}{0.5em}\setlength{\parindent}{0pt}}
\tcbset{before upper={\khoangcachdoan}}

\newtcolorbox{dinhnghia}[1][]{
  colback=box_bg, colframe=box_line, boxrule=0.6pt, arc=1pt,
  fonttitle=\bfseries, title={Định nghĩa #1}, breakable}
\newtcolorbox{dinhly}[1][]{
  colback=box_bg, colframe=cae_gold, boxrule=0.6pt, arc=1pt,
  fonttitle=\bfseries, title={Tính chất #1}, breakable}
\newtcolorbox{vidu}[1][]{
  colback=white, colframe=black!40, boxrule=0.4pt, arc=1pt,
  fonttitle=\bfseries, title={Ví dụ #1}, breakable}
\newtcolorbox{baitap}[1][]{
  colback=white, colframe=cae_blue!50, boxrule=0.4pt, arc=1pt,
  fonttitle=\bfseries, title={Câu #1}, breakable}

% Ô đáp án: viền xanh đậm hơn, dùng trong file đáp án
\newtcolorbox{dapan}[1][]{
  colback=box_bg, colframe=cae_blue, boxrule=0.6pt, arc=1pt,
  fonttitle=\bfseries, title={Câu #1}, breakable}

% ---------- TÊN CHỨNG MINH ----------
\renewcommand{\proofname}{\normalfont\bfseries Chứng minh}

% ---------- MÔN HỌC ----------
% Đổi được ở file .tex trước khi % =====================================================================
%  CAE SHANGHAI — GÓI LỆNH DÙNG CHUNG CHO TÀI LIỆU LUYỆN THI CSCA TOÁN
%  Biên dịch bằng XeLaTeX. Mọi file .tex trong buổi học \input file này.
% =====================================================================

% ---------- NGÔN NGỮ & PHÔNG CHỮ ----------
\usepackage{fontspec}
\usepackage{polyglossia}
\setmainlanguage{vietnamese}
\setotherlanguage{english}

% Chữ Latin (Anh + Việt): TeX Gyre Termes — bản clone Times New Roman,
% đi kèm TeX Live, hỗ trợ đầy đủ dấu tiếng Việt.
\setmainfont{TeX Gyre Termes}
\setsansfont{TeX Gyre Heros}

% Chữ Hán: ưu tiên Source Han Serif SC (思源宋体) — chuẩn học thuật Trung Quốc.
% Nếu máy chưa cài, tự động dùng Songti SC (宋体, có sẵn trên macOS).
% Cả hai đều thuộc họ 宋体 nên diện mạo tài liệu không đổi.
\IfFontExistsTF{Source Han Serif SC}{
  \newcommand{\cjkfamilyname}{Source Han Serif SC}
}{
  \IfFontExistsTF{Songti SC}{
    \newcommand{\cjkfamilyname}{Songti SC}
  }{
    \newcommand{\cjkfamilyname}{Noto Serif CJK SC}
  }
}

% xeCJK tự chuyển font khi gặp ký tự Hán, không cần bọc \cjk{} thủ công.
% Đây là điều kiện bắt buộc: font Latin (TeX Gyre Termes) không có glyph chữ Hán,
% nếu không khai báo xeCJK thì mọi chữ Hán viết trần sẽ bị nuốt khỏi bản in.
\usepackage{xeCJK}
\setCJKfamilyfont{caehan}{\cjkfamilyname}[Scale=0.95]
\setCJKmainfont{\cjkfamilyname}[Scale=0.95]
\setCJKsansfont{\cjkfamilyname}[Scale=0.95]
\setCJKmonofont{\cjkfamilyname}[Scale=0.95]

% \cjk{} ép dùng font Hán trong ngữ cảnh cần thiết (ví dụ trong tiêu đề in đậm
% hoặc trong bảng có định dạng riêng). Dùng \CJKfamily chứ không phải tên font,
% vì tên font chỉ là văn bản — sẽ in ra chữ thay vì đổi font.
\newcommand{\cjk}[1]{{\CJKfamily{caehan}#1}}

% ---------- TOÁN HỌC ----------
\usepackage{amsmath, amssymb, amsthm, mathtools}
\usepackage{unicode-math}
\setmathfont{TeX Gyre Termes Math}

% ---------- BỐ CỤC TRANG ----------
% Lề phải chừa vùng an toàn của letterhead:
%   banner trên chiếm 0–19mm  → top tối thiểu 2,2cm
%   banner dưới chiếm 281–297mm → bottom tối thiểu 2,2cm
% Lề trên 2,6cm (26mm) để tiêu đề đầu trang không đè banner (banner hết ở 19mm).
\usepackage[a4paper, top=2.6cm, bottom=2.4cm, left=2.5cm, right=2.0cm]{geometry}
\usepackage{fancyhdr}
\usepackage{titlesec}
\usepackage{enumitem}
\usepackage{graphicx}
\usepackage{booktabs, array, longtable, tabularx}
\usepackage{xcolor}
\usepackage[most]{tcolorbox}
\usepackage[colorlinks=true, linkcolor=cae_blue, citecolor=cae_blue, urlcolor=cae_blue]{hyperref}

% ---------- BẢNG MÀU THƯƠNG HIỆU ----------
\definecolor{cae_blue}{RGB}{0, 51, 102}
\definecolor{cae_gold}{RGB}{196, 149, 60}
\definecolor{box_bg}{RGB}{245, 247, 250}
\definecolor{box_line}{RGB}{0, 51, 102}

% ---------- TIÊU ĐỀ MỤC ----------
% Đệm đầu trang: banner letterhead chiếm 0–19mm, mà lề trên là 22mm.
% \section* đầu trang (sau \clearpage) sát lề nên phải chừa thêm, nếu không
% sẽ đè lên banner.
\titleformat{\section}
  {\normalfont\large\bfseries\sffamily\color{cae_blue}}{\thesection.}{0.6em}{}
\titleformat{\subsection}
  {\normalfont\normalsize\bfseries\sffamily\color{cae_blue!85!black}}{\thesubsection}{0.6em}{}
\titlespacing*{\section}{0pt}{14pt}{6pt}
\titlespacing*{\subsection}{0pt}{10pt}{4pt}

% Tiêu đề không đánh số ở đầu trang: \section* sau \clearpage nằm sát lề trên
% nên đè lên banner letterhead. Dùng \vspace* ngay TRƯỚC \section* không đủ,
% vì titlesec chèn khoảng cách riêng. Cách chắc chắn là đệm bằng \null.
\newcommand{\sectiondautrang}[1]{\section*{#1}}

% Dạng bài dùng \subsection*{} nên cần định dạng riêng, đậm và có gạch chân màu
\titleformat{\subsubsection}
  {\normalfont\normalsize\bfseries\sffamily\color{cae_blue!85!black}}{\thesubsubsection}{0.6em}{}

% ---------- LETTERHEAD (ảnh nền toàn trang) ----------
% Ảnh letterhead của CAE Shanghai: banner xanh trên/dưới kèm logo chìm giữa.
% Tỉ lệ ảnh 1414x2000 = 0.7070, khớp A4 (0.7071) nên chèn thẳng không méo.
% Đặt ở lớp nền dưới cùng để chữ không bị ảnh che.
% Dùng \AddToShipoutPictureBG (không có dấu *) để áp dụng cho MỌI trang.
\usepackage{eso-pic}
\AddToShipoutPictureBG{%
  \AtPageLowerLeft{%
    \includegraphics[width=\paperwidth,height=\paperheight]{images/letterhead.png}%
  }%
}

% ---------- ĐẦU / CHÂN TRANG ----------
% Letterhead đã có banner trên và dưới, nên bỏ dòng chạy để tránh chồng chữ.
% Số trang đặt trong vùng nội dung, sát trên banner dưới, lề phải.
\pagestyle{fancy}
\setlength{\headheight}{0pt}
\fancyhf{}
\renewcommand{\headrulewidth}{0pt}
\renewcommand{\footrulewidth}{0pt}
\fancyfoot[R]{\raisebox{0.9cm}{\small\color{cae_blue}\thepage}}

% ---------- MÔI TRƯỜNG NỘI DUNG ----------
% Trong tcolorbox, LaTeX đặt parindent = parskip = 0, nên các đoạn văn trong
% khung dính liền nhau. Đặt lại khoảng cách đoạn cho mọi khung.
\newcommand{\khoangcachdoan}{\setlength{\parskip}{0.5em}\setlength{\parindent}{0pt}}
\tcbset{before upper={\khoangcachdoan}}

\newtcolorbox{dinhnghia}[1][]{
  colback=box_bg, colframe=box_line, boxrule=0.6pt, arc=1pt,
  fonttitle=\bfseries, title={Định nghĩa #1}, breakable}
\newtcolorbox{dinhly}[1][]{
  colback=box_bg, colframe=cae_gold, boxrule=0.6pt, arc=1pt,
  fonttitle=\bfseries, title={Tính chất #1}, breakable}
\newtcolorbox{vidu}[1][]{
  colback=white, colframe=black!40, boxrule=0.4pt, arc=1pt,
  fonttitle=\bfseries, title={Ví dụ #1}, breakable}
\newtcolorbox{baitap}[1][]{
  colback=white, colframe=cae_blue!50, boxrule=0.4pt, arc=1pt,
  fonttitle=\bfseries, title={Câu #1}, breakable}

% Ô đáp án: viền xanh đậm hơn, dùng trong file đáp án
\newtcolorbox{dapan}[1][]{
  colback=box_bg, colframe=cae_blue, boxrule=0.6pt, arc=1pt,
  fonttitle=\bfseries, title={Câu #1}, breakable}

% ---------- TÊN CHỨNG MINH ----------
\renewcommand{\proofname}{\normalfont\bfseries Chứng minh}

% ---------- MÔN HỌC ----------
% Đổi được ở file .tex trước khi \input{cae-style} bằng \renewcommand{\monhoc}{...}
\providecommand{\monhoc}{Môn Toán}

% ---------- TRANG BÌA GỌN ----------
% #1 = loại tài liệu (TÀI LIỆU CHUẨN BỊ, TÀI LIỆU BUỔI HỌC, ...)
% #2 = số buổi
% #3 = tên buổi tiếng Việt
% #4 = tên buổi tiếng Hán
% #5 = module và tỉ trọng
\newcommand{\trangbia}[5]{%
  \begin{center}
    \vspace*{0.6cm}
    {\large\bfseries\color{cae_blue} CAE SHANGHAI}\\[0.15cm]
    {\footnotesize Tài liệu luyện thi CSCA · \monhoc}\\[0.5cm]
    {\color{cae_gold}\rule{0.55\linewidth}{0.8pt}}\\[0.35cm]
    {\small\bfseries\color{cae_blue!80!black} #1}\\[0.35cm]
    {\Large\bfseries\color{cae_blue} BUỔI #2 — #3}\\[0.2cm]
    {\large\color{cae_blue!85!black} \cjk{#4}}\\[0.25cm]
    \ifx\relax#5\relax\else{\small #5}\\[0.35cm]\fi
    {\color{cae_gold}\rule{0.55\linewidth}{0.8pt}}\\[0.3cm]
    {\footnotesize Biên soạn: Hoàng Dương, CAE Shanghai \quad·\quad \today}
  \end{center}
  \vspace{0.3cm}
}

% ---------- TRANG BÌA CHO ĐỀ KIỂM TRA ----------
% Khác \trangbia ở chỗ không có chữ "BUỔI N" — dùng cho đề thi, không phải bài giảng.
% #1 = loại tài liệu, #2 = tên kỳ thi (Việt), #3 = tên kỳ thi (Hán), #4 = thông số
\newcommand{\biade}[4]{%
  \begin{center}
    \vspace*{0.6cm}
    {\large\bfseries\color{cae_blue} CAE SHANGHAI}\\[0.15cm]
    {\footnotesize Tài liệu luyện thi CSCA · \monhoc}\\[0.5cm]
    {\color{cae_gold}\rule{0.55\linewidth}{0.8pt}}\\[0.35cm]
    {\small\bfseries\color{cae_blue!80!black} #1}\\[0.35cm]
    {\Large\bfseries\color{cae_blue} #2}\\[0.2cm]
    {\large\color{cae_blue!85!black} \cjk{#3}}\\[0.25cm]
    {\small #4}\\[0.35cm]
    {\color{cae_gold}\rule{0.55\linewidth}{0.8pt}}\\[0.3cm]
    {\footnotesize Biên soạn: Hoàng Dương, CAE Shanghai \quad·\quad \today}
  \end{center}
  \vspace{0.3cm}
}

% ---------- TIỆN ÍCH ----------
% Nhãn từ khóa in đậm ở đầu đoạn, dùng cho Từ vựng / Từ khóa nhận dạng / Chú ý
% Khoảng cách đoạn do khung tự đặt (xem \khoangcachdoan ở trên),
% ở đây chỉ cần tách đoạn và bỏ thụt đầu dòng.
\newcommand{\nhan}[1]{\par\noindent\textbf{#1.}\ \ignorespaces}

% Bảng gọn, không đánh số
\newcommand{\bang}[1]{%
  \begin{center}#1\end{center}}

% Đánh dấu ô trống cho học sinh điền
\newcommand{\chotrong}{\makebox[3.2cm]{\dotfill}}
 bằng \renewcommand{\monhoc}{...}
\providecommand{\monhoc}{Môn Toán}

% ---------- TRANG BÌA GỌN ----------
% #1 = loại tài liệu (TÀI LIỆU CHUẨN BỊ, TÀI LIỆU BUỔI HỌC, ...)
% #2 = số buổi
% #3 = tên buổi tiếng Việt
% #4 = tên buổi tiếng Hán
% #5 = module và tỉ trọng
\newcommand{\trangbia}[5]{%
  \begin{center}
    \vspace*{0.6cm}
    {\large\bfseries\color{cae_blue} CAE SHANGHAI}\\[0.15cm]
    {\footnotesize Tài liệu luyện thi CSCA · \monhoc}\\[0.5cm]
    {\color{cae_gold}\rule{0.55\linewidth}{0.8pt}}\\[0.35cm]
    {\small\bfseries\color{cae_blue!80!black} #1}\\[0.35cm]
    {\Large\bfseries\color{cae_blue} BUỔI #2 — #3}\\[0.2cm]
    {\large\color{cae_blue!85!black} \cjk{#4}}\\[0.25cm]
    \ifx\relax#5\relax\else{\small #5}\\[0.35cm]\fi
    {\color{cae_gold}\rule{0.55\linewidth}{0.8pt}}\\[0.3cm]
    {\footnotesize Biên soạn: Hoàng Dương, CAE Shanghai \quad·\quad \today}
  \end{center}
  \vspace{0.3cm}
}

% ---------- TRANG BÌA CHO ĐỀ KIỂM TRA ----------
% Khác \trangbia ở chỗ không có chữ "BUỔI N" — dùng cho đề thi, không phải bài giảng.
% #1 = loại tài liệu, #2 = tên kỳ thi (Việt), #3 = tên kỳ thi (Hán), #4 = thông số
\newcommand{\biade}[4]{%
  \begin{center}
    \vspace*{0.6cm}
    {\large\bfseries\color{cae_blue} CAE SHANGHAI}\\[0.15cm]
    {\footnotesize Tài liệu luyện thi CSCA · \monhoc}\\[0.5cm]
    {\color{cae_gold}\rule{0.55\linewidth}{0.8pt}}\\[0.35cm]
    {\small\bfseries\color{cae_blue!80!black} #1}\\[0.35cm]
    {\Large\bfseries\color{cae_blue} #2}\\[0.2cm]
    {\large\color{cae_blue!85!black} \cjk{#3}}\\[0.25cm]
    {\small #4}\\[0.35cm]
    {\color{cae_gold}\rule{0.55\linewidth}{0.8pt}}\\[0.3cm]
    {\footnotesize Biên soạn: Hoàng Dương, CAE Shanghai \quad·\quad \today}
  \end{center}
  \vspace{0.3cm}
}

% ---------- TIỆN ÍCH ----------
% Nhãn từ khóa in đậm ở đầu đoạn, dùng cho Từ vựng / Từ khóa nhận dạng / Chú ý
% Khoảng cách đoạn do khung tự đặt (xem \khoangcachdoan ở trên),
% ở đây chỉ cần tách đoạn và bỏ thụt đầu dòng.
\newcommand{\nhan}[1]{\par\noindent\textbf{#1.}\ \ignorespaces}

% Bảng gọn, không đánh số
\newcommand{\bang}[1]{%
  \begin{center}#1\end{center}}

% Đánh dấu ô trống cho học sinh điền
\newcommand{\chotrong}{\makebox[3.2cm]{\dotfill}}
 bằng \renewcommand{\monhoc}{...}
\providecommand{\monhoc}{Môn Toán}

% ---------- TRANG BÌA GỌN ----------
% #1 = loại tài liệu (TÀI LIỆU CHUẨN BỊ, TÀI LIỆU BUỔI HỌC, ...)
% #2 = số buổi
% #3 = tên buổi tiếng Việt
% #4 = tên buổi tiếng Hán
% #5 = module và tỉ trọng
\newcommand{\trangbia}[5]{%
  \begin{center}
    \vspace*{0.6cm}
    {\large\bfseries\color{cae_blue} CAE SHANGHAI}\\[0.15cm]
    {\footnotesize Tài liệu luyện thi CSCA · \monhoc}\\[0.5cm]
    {\color{cae_gold}\rule{0.55\linewidth}{0.8pt}}\\[0.35cm]
    {\small\bfseries\color{cae_blue!80!black} #1}\\[0.35cm]
    {\Large\bfseries\color{cae_blue} BUỔI #2 — #3}\\[0.2cm]
    {\large\color{cae_blue!85!black} \cjk{#4}}\\[0.25cm]
    \ifx\relax#5\relax\else{\small #5}\\[0.35cm]\fi
    {\color{cae_gold}\rule{0.55\linewidth}{0.8pt}}\\[0.3cm]
    {\footnotesize Biên soạn: Hoàng Dương, CAE Shanghai \quad·\quad \today}
  \end{center}
  \vspace{0.3cm}
}

% ---------- TRANG BÌA CHO ĐỀ KIỂM TRA ----------
% Khác \trangbia ở chỗ không có chữ "BUỔI N" — dùng cho đề thi, không phải bài giảng.
% #1 = loại tài liệu, #2 = tên kỳ thi (Việt), #3 = tên kỳ thi (Hán), #4 = thông số
\newcommand{\biade}[4]{%
  \begin{center}
    \vspace*{0.6cm}
    {\large\bfseries\color{cae_blue} CAE SHANGHAI}\\[0.15cm]
    {\footnotesize Tài liệu luyện thi CSCA · \monhoc}\\[0.5cm]
    {\color{cae_gold}\rule{0.55\linewidth}{0.8pt}}\\[0.35cm]
    {\small\bfseries\color{cae_blue!80!black} #1}\\[0.35cm]
    {\Large\bfseries\color{cae_blue} #2}\\[0.2cm]
    {\large\color{cae_blue!85!black} \cjk{#3}}\\[0.25cm]
    {\small #4}\\[0.35cm]
    {\color{cae_gold}\rule{0.55\linewidth}{0.8pt}}\\[0.3cm]
    {\footnotesize Biên soạn: Hoàng Dương, CAE Shanghai \quad·\quad \today}
  \end{center}
  \vspace{0.3cm}
}

% ---------- TIỆN ÍCH ----------
% Nhãn từ khóa in đậm ở đầu đoạn, dùng cho Từ vựng / Từ khóa nhận dạng / Chú ý
% Khoảng cách đoạn do khung tự đặt (xem \khoangcachdoan ở trên),
% ở đây chỉ cần tách đoạn và bỏ thụt đầu dòng.
\newcommand{\nhan}[1]{\par\noindent\textbf{#1.}\ \ignorespaces}

% Bảng gọn, không đánh số
\newcommand{\bang}[1]{%
  \begin{center}#1\end{center}}

% Đánh dấu ô trống cho học sinh điền
\newcommand{\chotrong}{\makebox[3.2cm]{\dotfill}}


\begin{document}

\trangbia{ĐÁP ÁN ĐỀ CHỮA}{19}{Toàn bộ chương trình}{综合模拟}{Đáp án và lời giải}

\section{BẢNG ĐÁP ÁN}

\bang{
\begin{tabular}{@{}lcccccccccccc@{}}
\toprule
\textbf{Câu} & \textbf{1} & \textbf{2} & \textbf{3} & \textbf{4} & \textbf{5} & \textbf{6} & \textbf{7} & \textbf{8} & \textbf{9} & \textbf{10} & \textbf{11} & \textbf{12} \\
\midrule
Đáp án & B & A & A & B & B & A & B & B & D & D & A & D \\
\bottomrule
\end{tabular}}

\bang{
\begin{tabular}{@{}lcccccccccccc@{}}
\toprule
\textbf{Câu} & \textbf{13} & \textbf{14} & \textbf{15} & \textbf{16} & \textbf{17} & \textbf{18} & \textbf{19} & \textbf{20} & \textbf{21} & \textbf{22} & \textbf{23} & \textbf{24} \\
\midrule
Đáp án & B & D & B & C & C & D & D & C & C & D & C & A \\
\bottomrule
\end{tabular}}

\bang{
\begin{tabular}{@{}lcccccccccccc@{}}
\toprule
\textbf{Câu} & \textbf{25} & \textbf{26} & \textbf{27} & \textbf{28} & \textbf{29} & \textbf{30} & \textbf{31} & \textbf{32} & \textbf{33} & \textbf{34} & \textbf{35} & \textbf{36} \\
\midrule
Đáp án & A & C & A & D & D & A & C & A & C & D & C & A \\
\bottomrule
\end{tabular}}

\bang{
\begin{tabular}{@{}lcccccccccccc@{}}
\toprule
\textbf{Câu} & \textbf{37} & \textbf{38} & \textbf{39} & \textbf{40} & \textbf{41} & \textbf{42} & \textbf{43} & \textbf{44} & \textbf{45} & \textbf{46} & \textbf{47} & \textbf{48} \\
\midrule
Đáp án & B & B & C & B & C & B & A & A & D & B & C & D \\
\bottomrule
\end{tabular}}

\section{PHẦN A — LƯỢNG GIÁC}

\begin{dapan}[1 — Đáp án B]

\textbf{Giải.} $\alpha$ thuộc góc phần tư thứ ba nên $\cos\alpha < 0$. Áp dụng $\sin^2\alpha + \cos^2\alpha = 1$:

$$\cos\alpha = -\sqrt{1 - \left(\dfrac{4}{5}\right)^2} = -\sqrt{\dfrac{9}{25}} = -\dfrac{3}{5}$$


\end{dapan}
\vspace{0.3cm}
\begin{dapan}[2 — Đáp án A]

\textbf{Giải.} Công thức góc bù: $\sin 150^\circ = \sin(180^\circ - 30^\circ) = \sin 30^\circ = \dfrac{1}{2}$.


\end{dapan}
\vspace{0.3cm}
\begin{dapan}[3 — Đáp án A]

\textbf{Giải.} Chu kì của $y = \sin(\omega x)$ là $T = \dfrac{2\pi}{\lvert \omega \rvert}$. Với $\omega = 3$:

$$T = \dfrac{2\pi}{3}$$


\end{dapan}
\vspace{0.3cm}
\begin{dapan}[4 — Đáp án B]

\textbf{Giải.} Vì $-1 \le \sin x \le 1$, biểu thức lớn nhất khi $\sin x$ nhỏ nhất:

$$y\_{\max} = 3 - 2 \cdot (-1) = 5$$


\end{dapan}
\vspace{0.3cm}
\begin{dapan}[5 — Đáp án B]

\textbf{Giải.} Chia cả tử và mẫu cho $\cos\alpha$ để quy về $\tan\alpha$:

$$\dfrac{\sin\alpha + \cos\alpha}{2\sin\alpha - \cos\alpha} = \dfrac{\tan\alpha + 1}{2\tan\alpha - 1} = \dfrac{3 + 1}{6 - 1} = \dfrac{4}{5}$$


\end{dapan}
\vspace{0.3cm}
\begin{dapan}[6 — Đáp án A]

\textbf{Giải.} Dạng $a\sin x + b\cos x$ có giá trị lớn nhất $\sqrt{a^2 + b^2}$. Với $a = b = 1$:

$$\sqrt{1 + 1} = \sqrt{2}$$


\end{dapan}
\vspace{0.3cm}
\begin{dapan}[7 — Đáp án B]

\textbf{Giải.} Góc phần tư thứ ba cho $\sin\alpha < 0$:

$$\sin\alpha = -\sqrt{1 - \dfrac{9}{25}} = -\dfrac{4}{5}, \qquad \tan\alpha = \dfrac{\sin\alpha}{\cos\alpha} = \dfrac{-4/5}{-3/5} = \dfrac{4}{3}$$


\end{dapan}
\vspace{0.3cm}
\begin{dapan}[8 — Đáp án B]

\textbf{Giải.} Bình phương đẳng thức đã cho:

$$(\sin\alpha + \cos\alpha)^2 = 1 + \sin 2\alpha = \dfrac{1}{4} \implies \sin 2\alpha = \dfrac{1}{4} - 1 = -\dfrac{3}{4}$$


\end{dapan}
\vspace{0.3cm}
\begin{dapan}[9 — Đáp án D]

\textbf{Giải.} Từ $\sin\alpha = \dfrac{3}{5}$ và $\alpha$ ở góc phần tư thứ nhất, $\cos\alpha = \dfrac{4}{5}$. Công thức nhân đôi:

$$\sin 2\alpha = 2\sin\alpha\cos\alpha = 2 \cdot \dfrac{3}{5} \cdot \dfrac{4}{5} = \dfrac{24}{25}$$


\end{dapan}
\vspace{0.3cm}
\begin{dapan}[10 — Đáp án D]

\textbf{Giải.} Dùng công thức hạ bậc ở tử:

$$1 - \cos 2\alpha = 2\sin^2\alpha, \qquad \sin 2\alpha = 2\sin\alpha\cos\alpha$$

$$\dfrac{1 - \cos 2\alpha}{\sin 2\alpha} = \dfrac{2\sin^2\alpha}{2\sin\alpha\cos\alpha} = \dfrac{\sin\alpha}{\cos\alpha} = \tan\alpha$$

\section{PHẦN B — HÌNH HỌC GIẢI TÍCH}


\end{dapan}
\vspace{0.3cm}
\begin{dapan}[11 — Đáp án A]

\textbf{Giải.} Trục $Ox$ ứng với $y = 0$. Thay vào: $2x + 6 = 0$, được $x = -3$.


\end{dapan}
\vspace{0.3cm}
\begin{dapan}[12 — Đáp án D]

\textbf{Giải.} Đường song song có cùng hệ số góc $2$, nên có dạng $y = 2x + b$. Thay điểm $(1, 2)$: $2 = 2 + b$, được $b = 0$. Vậy $y = 2x$.


\end{dapan}
\vspace{0.3cm}
\begin{dapan}[13 — Đáp án B]

\textbf{Giải.} Tâm $I(2, -3)$. Bán kính:

$$R = \sqrt{2^2 + (-3)^2 - (-4)} = \sqrt{4 + 9 + 4} = \sqrt{17}$$


\end{dapan}
\vspace{0.3cm}
\begin{dapan}[14 — Đáp án D]

\textbf{Giải.} Công thức khoảng cách từ điểm đến đường thẳng:

$$d = \dfrac{\lvert 3 \cdot 4 - 4 \cdot 1 + 2 \rvert}{\sqrt{3^2 + (-4)^2}} = \dfrac{\lvert 10 \rvert}{5} = 2$$


\end{dapan}
\vspace{0.3cm}
\begin{dapan}[15 — Đáp án B]

\textbf{Giải.} Dạng $y^2 = 2px$ với $2p = 8$, tức $p = 4$. Parabol mở sang phải, đường chuẩn là $x = -\dfrac{p}{2} = -2$.


\end{dapan}
\vspace{0.3cm}
\begin{dapan}[16 — Đáp án C]

\textbf{Giải.} Với $\dfrac{x^2}{25} + \dfrac{y^2}{9} = 1$: $a = 5$, $b = 3$, nên $c = \sqrt{25 - 9} = 4$. Khoảng cách giữa hai tiêu điểm là $2c = 8$.


\end{dapan}
\vspace{0.3cm}
\begin{dapan}[17 — Đáp án C]

\textbf{Giải.} Với $\dfrac{x^2}{9} - \dfrac{y^2}{16} = 1$: $a = 3$, $b = 4$, $c = \sqrt{9 + 16} = 5$. Tâm sai $e = \dfrac{c}{a} = \dfrac{5}{3}$.


\end{dapan}
\vspace{0.3cm}
\begin{dapan}[18 — Đáp án D]

\textbf{Giải.} Tâm đường tròn $x^2 + y^2 + Dx + Ey + F = 0$ là $\left(-\dfrac{D}{2}, -\dfrac{E}{2}\right)$. Với $D = -2$, $E = 4$, tâm là $(1, -2)$.


\end{dapan}
\vspace{0.3cm}
\begin{dapan}[19 — Đáp án D]

\textbf{Giải.} Đường $4x + 3y - 5 = 0$ có hệ số góc $-\dfrac{4}{3}$, nên đường vuông góc có hệ số góc $\dfrac{3}{4}$.

$$y - 1 = \dfrac{3}{4}(x - 2) \iff 3x - 4y - 2 = 0$$

\section{PHẦN C — DÃY SỐ}


\end{dapan}
\vspace{0.3cm}
\begin{dapan}[20 — Đáp án C]

\textbf{Giải.} $a_4 = a_1 + 3d = 2 + 9 = 11$.


\end{dapan}
\vspace{0.3cm}
\begin{dapan}[21 — Đáp án C]

\textbf{Giải.} Tổng $n$ số hạng đầu của cấp số cộng: $S_n = \dfrac{n}{2}\left[2a_1 + (n-1)d\right]$.

$$S\_5 = \dfrac{5}{2}\left[2 \cdot 10 + 4 \cdot (-2)\right] = \dfrac{5}{2} \cdot 12 = 30$$


\end{dapan}
\vspace{0.3cm}
\begin{dapan}[22 — Đáp án D]

\textbf{Giải.} $a_7 - a_3 = 4d = 12$ nên $d = 3$. Vậy $a_1 = a_3 - 2d = 7 - 6 = 1$.


\end{dapan}
\vspace{0.3cm}
\begin{dapan}[23 — Đáp án C]

\textbf{Giải.} $a_4 = a_1 q^3 = 3 \cdot 2^3 = 24$.


\end{dapan}
\vspace{0.3cm}
\begin{dapan}[24 — Đáp án A]

\textbf{Giải.} $a_4 = a_2 q^2$ nên $24 = 6q^2$, được $q^2 = 4$. Vì các số hạng dương nên $q = 2$. Vậy $a_1 = \dfrac{a_2}{q} = \dfrac{6}{2} = 3$.


\end{dapan}
\vspace{0.3cm}
\begin{dapan}[25 — Đáp án A]

\textbf{Giải.} Tính lần lượt theo hệ thức truy hồi:

$$a\_2 = 2 \cdot 1 + 1 = 3, \quad a\_3 = 2 \cdot 3 + 1 = 7, \quad a\_4 = 2 \cdot 7 + 1 = 15$$

\section{PHẦN D — HÀM SỐ}


\end{dapan}
\vspace{0.3cm}
\begin{dapan}[26 — Đáp án C]

\textbf{Giải.} Căn thức cần $2x - 6 \ge 0$, tức $x \ge 3$. Mẫu thức cần $x \ne 5$. Kết hợp: $[3, 5) \cup (5, +\infty)$.


\end{dapan}
\vspace{0.3cm}
\begin{dapan}[27 — Đáp án A]

\textbf{Giải.} Trục đối xứng của parabol $y = ax^2 + bx + c$ là $x = -\dfrac{b}{2a}$. Với $a = 1$, $b = -6$:

$$x = \dfrac{6}{2} = 3$$


\end{dapan}
\vspace{0.3cm}
\begin{dapan}[28 — Đáp án D]

\textbf{Giải.} $f'(x) = 3x^2 - 6x = 3x(x - 2)$. Trên $[0, 2]$ đạo hàm âm trong $(0, 2)$, nên hàm nghịch biến và giá trị nhỏ nhất tại $x = 2$:

$$f(2) = 8 - 12 + 1 = -3$$


\end{dapan}
\vspace{0.3cm}
\begin{dapan}[29 — Đáp án D]

\textbf{Giải.} $f(2) = 3 \cdot 2 - 2 = 4$.


\end{dapan}
\vspace{0.3cm}
\begin{dapan}[30 — Đáp án A]

\textbf{Giải.} Hàm logarit cần biểu thức trong dấu logarit dương: $x - 1 > 0$, tức $x > 1$. Tập xác định là $(1, +\infty)$.


\end{dapan}
\vspace{0.3cm}
\begin{dapan}[31 — Đáp án C]

\textbf{Giải.} Xét $y = \dfrac{2x}{x^2 + 1}$. Đặt $y = k$, phương trình $kx^2 - 2x + k = 0$ phải có nghiệm, nên $\Delta' = 1 - k^2 \ge 0$, tức $-1 \le k \le 1$. Giá trị lớn nhất là $1$ (đạt tại $x = 1$).

\section{PHẦN E — TẬP HỢP VÀ BẤT ĐẲNG THỨC}


\end{dapan}
\vspace{0.3cm}
\begin{dapan}[32 — Đáp án A]

\textbf{Giải.} $x^2 - 9 < 0$ cho $A = (-3, 3)$. Giao với $B = [0, +\infty)$ là $[0, 3)$.


\end{dapan}
\vspace{0.3cm}
\begin{dapan}[33 — Đáp án C]

\textbf{Giải.} Phần chung của hai tập hợp gồm các phần tử có mặt ở cả hai: $\{2, 3\}$.


\end{dapan}
\vspace{0.3cm}
\begin{dapan}[34 — Đáp án D]

\textbf{Giải.} $\lvert 2x - 1 \rvert < 3 \iff -3 < 2x - 1 < 3$. Cộng $1$: $-2 < 2x < 4$. Chia $2$: $-1 < x < 2$.


\end{dapan}
\vspace{0.3cm}
\begin{dapan}[35 — Đáp án C]

\textbf{Giải.} $x^2 - 5x + 6 = (x - 2)(x - 3)$. Vì $a = 1 > 0$, biểu thức âm trong khoảng giữa hai nghiệm: $2 < x < 3$.


\end{dapan}
\vspace{0.3cm}
\begin{dapan}[36 — Đáp án A]

\textbf{Giải.} Với $x > 0$, áp dụng bất đẳng thức Cauchy:

$$x + \dfrac{4}{x} \ge 2\sqrt{x \cdot \dfrac{4}{x}} = 4$$

Dấu bằng khi $x = 2$.

\section{PHẦN F — MŨ VÀ LOGARIT}


\end{dapan}
\vspace{0.3cm}
\begin{dapan}[37 — Đáp án B]

\textbf{Giải.} $32 = 2^5$ nên $\log_2 32 = 5$.


\end{dapan}
\vspace{0.3cm}
\begin{dapan}[38 — Đáp án B]

\textbf{Giải.} Dùng $a^{\log_a b} = b$ với $a = 3$, $b = 5$: $3^{\log_3 5} = 5$.


\end{dapan}
\vspace{0.3cm}
\begin{dapan}[39 — Đáp án C]

\textbf{Giải.} Lấy logarit cơ số $2$ hai vế: $x = \log_2 5$.


\end{dapan}
\vspace{0.3cm}
\begin{dapan}[40 — Đáp án B]

\textbf{Giải.} Viết $4 = \left(\dfrac{1}{2}\right)^{-2}$. Vì cơ số $\dfrac{1}{2} < 1$, đổi chiều bất đẳng thức:

$$\left(\dfrac{1}{2}\right)^x < \left(\dfrac{1}{2}\right)^{-2} \iff x > -2$$

\section{PHẦN G — VECTOR VÀ SỐ PHỨC}


\end{dapan}
\vspace{0.3cm}
\begin{dapan}[41 — Đáp án C]

\textbf{Giải.} Tích vô hướng: $a \cdot b = 2 \cdot 1 + (-1) \cdot 3 = 2 - 3 = -1$.


\end{dapan}
\vspace{0.3cm}
\begin{dapan}[42 — Đáp án B]

\textbf{Giải.} Cộng phần thực với phần thực, phần ảo với phần ảo:

$$(2 - i) + (1 + 3i) = (2 + 1) + (-1 + 3)i = 3 + 2i$$


\end{dapan}
\vspace{0.3cm}
\begin{dapan}[43 — Đáp án A]

\textbf{Giải.} $\lvert z \rvert = \sqrt{3^2 + (-4)^2} = \sqrt{25} = 5$.

\section{PHẦN H — XÁC SUẤT VÀ THỐNG KÊ}


\end{dapan}
\vspace{0.3cm}
\begin{dapan}[44 — Đáp án A]

\textbf{Giải.} Tổng số bóng là $5$, số bóng đỏ là $2$. Xác suất là $\dfrac{2}{5}$.


\end{dapan}
\vspace{0.3cm}
\begin{dapan}[45 — Đáp án D]

\textbf{Giải.} Sắp xếp dữ liệu: $2, 4, 4, 6, 9$. Có $5$ giá trị nên trung vị là giá trị thứ ba, bằng $4$.

\section{PHẦN I — HÌNH HỌC KHÔNG GIAN}


\end{dapan}
\vspace{0.3cm}
\begin{dapan}[46 — Đáp án B]

\textbf{Giải.} So với dạng $(x-a)^2 + (y-b)^2 + (z-c)^2 = R^2$, tâm là $(2, -1, 3)$.


\end{dapan}
\vspace{0.3cm}
\begin{dapan}[47 — Đáp án C]

\textbf{Giải.} Khoảng cách từ $(1, 1, 1)$ đến gốc tọa độ:

$$d = \sqrt{1^2 + 1^2 + 1^2} = \sqrt{3}$$


\end{dapan}
\vspace{0.3cm}
\begin{dapan}[48 — Đáp án D]

\textbf{Giải.} Diện tích mặt cầu $S = 4\pi R^2 = 4\pi \cdot 4 = 16\pi$.

\emph{Tài liệu nội bộ, CAE SHANGHAI.}
\end{dapan}
\vspace{0.3cm}

\vspace{0.8cm}
\begin{center}
{\footnotesize\color{cae_blue} Tài liệu nội bộ, CAE SHANGHAI}
\end{center}

\end{document}
